% Generated by tools/edn_latex.py from reports/edn.md.
% Do not edit: change reports/edn.md and run `make edn`.
\ednentry{
  id = {EDN-51},
  title = {Mean fill plus a per-feature missingness indicator for the Ridge numerics, which is not imputation},
  date = {2026-09-30},
  milestone = {M3: Quality Assurance},
  topic = {Feature Engineering, Modelling Approach},
  participants = {@lukas2510},
  decision = {The \texttt{b1} Ridge variant's numeric branch is the union of two paths: a mean-filled, standardised copy of every numeric feature, and \texttt{MissingIndicator(features=\allowbreak{}\textquotedbl{}all\textquotedbl{})}, which emits one indicator column per feature whether or not that feature was ever missing at fit time. No other estimator fills anything, and no artefact of the pipeline is imputed: the fill exists only inside B1's fitted pipeline.},
  rationale = {\emph{Alternatives considered:}
\begin{itemize}[leftmargin=*, topsep=2pt]
\item \textbf{Option A (chosen): mean fill plus an indicator for every feature.}
\newline Pros: information-preserving for a linear model. For a row where feature \texttt{j} is absent the contribution is \texttt{beta\_\allowbreak{}j * mean\_\allowbreak{}j + gamma\_\allowbreak{}j}, and \texttt{gamma\_\allowbreak{}j} is free to absorb whatever the absence is worth, so the fill value is a numerically neutral placeholder rather than a guess at the value. Nothing about the missingness is destroyed, which is what EDN-15 is protecting. \texttt{features=\allowbreak{}\textquotedbl{}all\textquotedbl{}} also makes SC-06 honest: the criterion masks a field and measures what it costs, and it can only do that if masking a field the training rows happened to fill still produces an indicator.
\newline Cons: a reader who sees \texttt{SimpleImputer} in the pipeline will conclude the project imputes, which is exactly the misreading this entry exists to prevent. It also doubles the numeric column count.
\item \textbf{Option B: \texttt{SimpleImputer} alone, or \texttt{SimpleImputer(add\_\allowbreak{}indicator=\allowbreak{}True)} at its default.}
\newline Pros: less to explain, fewer columns.
\newline Cons: this is the option that silently breaks SC-06, and it is the default, which is why it is worth an entry. scikit-learn emits an indicator only for features that were missing \textbf{at fit time}, so a column complete in the training split is mean-filled with no indicator the moment the criterion masks it -- the model then reports the training mean for that field, the prediction barely moves, and SC-06 measures the model's own imputation rather than the cost of the missing input. Without the indicator at all, an absent value is indistinguishable from an average one, which is imputation in the sense EDN-15 rejects.
\item \textbf{Option C: drop the rows or the columns with missing values.}
\newline Pros: no fill anywhere, so no explaining to do.
\newline Cons: dropping rows would fit B1 on a different and much smaller population than the other three variants, so the ladder would stop being a comparison. Dropping columns would throw away \texttt{nr\_\allowbreak{}prev\_\allowbreak{}owners} and \texttt{gears}, which are missing in a third and a quarter of the training rows respectively, and the missingness itself carries signal (EDN-15).
\item \textbf{Option D: give up on B1 taking numerics with gaps and use a model that handles them natively.}
\newline Pros: no encoding at all.
\newline Cons: that model is LightGBM, which is ladder steps 3 and 4. B1's value is that it is a \emph{linear} reference.
\end{itemize}
\par
\emph{Why:} EDN-15's rule is that a missing value is never imputed, because the missingness itself carries signal. Ridge cannot consume NaN, so B1 either encodes the gaps or does not exist. The construction above is the encoding under which nothing is lost: the pair \texttt{(filled value, indicator)} is a bijection with \texttt{(value, present)} for a present value and carries the absence explicitly for an absent one, so the linear model has exactly the information the tree models get from a NaN split. The fill is a placeholder chosen to be numerically harmless after standardisation, not an estimate of the missing value, and the entry exists because those two look identical in the code.
\par
\texttt{features=\allowbreak{}\textquotedbl{}all\textquotedbl{}} is the part a test has to hold, because it is invisible from the outside. It is also invisible on the synthetic fixture: every numeric column there has a missing value in the training split, so the default emits the same set of indicators and a test written against the fixture cannot tell the two apart. That was found by breaking the implementation on purpose and watching the test suite pass. The test now fills one column completely before fitting and fails with the default.
\par
What that test constructs is \textbf{not} the situation on the real snapshot, and the entry says so rather than borrowing the stronger claim. Of the 8 numeric features of the \texttt{basic} set -- the only set \texttt{b1} uses -- \textbf{none} is complete in the real training split: the emptiest is \texttt{age\_\allowbreak{}years} with 1 gap in 60,378 rows and the fullest \texttt{nr\_\allowbreak{}prev\_\allowbreak{}owners} with 22,697. So scikit-learn's default would emit the same eight indicators here, and \texttt{features=\allowbreak{}\textquotedbl{}all\textquotedbl{}} changes nothing in the shipped configuration. It is insurance, not a measured fix, and what it insures against is real: 136 of the \texttt{extended} set's 147 numeric features are complete at fit time, so a Ridge on \texttt{extended} -- which the ladder may well want, since B1 is the interpretable reference for the winning feature set -- would lose 136 of its 147 indicators to the default, and SC-06 would measure the model's own imputation for every one of them. The guard is kept because it costs nothing and because the configuration it protects is one change away.
\par
The other half of the same discipline is what the estimator does \emph{not} do. \texttt{SimpleImputer(keep\_\allowbreak{}empty\_\allowbreak{}features=\allowbreak{}True)} keeps an all-missing column in place rather than dropping it, so the matrix shape does not depend on which columns the training rows happened to fill; and nothing outside B1's pipeline fills anything, so no artefact on disk and no other variant carries a filled value.},
  ai = {Alternative generation, Alternative assessment, Recommendation, Solution generation},
  response = {Accepted},
  assessment = {AI identified that the obvious construction would make SC-06 measure its own imputation, which is a subtle failure that would have produced a passing criterion and a wrong conclusion, and it named \texttt{features=\allowbreak{}\textquotedbl{}all\textquotedbl{}} as the fix and flagged the whole construction as an EDN candidate because of how it reads. It also proposed the test for it. The test it proposed was then shown, by mutating the implementation, to pass either way on this fixture, and it was replaced with one that fills a column first; that correction came from the mutation exercise rather than from the design.},
  aievidence = {Claude Code session on 2026-09-30 implementing issue \#37: the mutation battery run before the pull request recorded \texttt{MissingIndicator(features=\allowbreak{}\textquotedbl{}all\textquotedbl{})} to \texttt{MissingIndicator()} as a surviving mutation, which is what produced the replacement test.},
  evidence = {\href{https://github.com/mlops-2627q1-mds-upc/MLOPS_Recommenditos/blob/main/recommenditos/modeling/model.py}{\texttt{recommenditos/\allowbreak{}modeling/\allowbreak{}model.\allowbreak{}py}}, \texttt{RidgeModel} and \texttt{\_\allowbreak{}ridge\_\allowbreak{}pipeline}; \texttt{tests/\allowbreak{}test\_\allowbreak{}model.\allowbreak{}py::test\_\allowbreak{}the\_\allowbreak{}ridge\_\allowbreak{}emits\_\allowbreak{}a\_\allowbreak{}missing\_\allowbreak{}indicator\_\allowbreak{}for\_\allowbreak{}every\_\allowbreak{}numeric\_\allowbreak{}feature} and \texttt{::test\_\allowbreak{}a\_\allowbreak{}numeric\_\allowbreak{}feature\_\allowbreak{}complete\_\allowbreak{}in\_\allowbreak{}training\_\allowbreak{}still\_\allowbreak{}gets\_\allowbreak{}an\_\allowbreak{}indicator}; \href{https://github.com/mlops-2627q1-mds-upc/MLOPS_Recommenditos/blob/main/docs/docs/model-card.md}{model card}, Training Procedure, Training; \hyperref[EDN-15]{EDN-15}; \href{https://github.com/mlops-2627q1-mds-upc/MLOPS_Recommenditos/issues/37}{issue \#37}.},
}
